\begin{proof}
The mass and mean assertions are immediate. The impulse Hankel
$H_{\delta_{w_L}}$ is a partial permutation of norm one. Subtracting
the rank-one operator $(1-\delta_L)H_{\delta_\varepsilon}$ therefore
gives $\sigma_2(H_{p_L})\le\delta_L$.

For any $q\in\mathcal Q_N^+$, pad its realization with unused states
if necessary, and let $T=\sum_aT_a$. Since $T\mathbf1\le\mathbf1$,
also $T^m\mathbf1\le\mathbf1$. For $B\in\Sigma^m$ put
$\psi(B)=\|T_B\|_{\infty\to\infty}
=\max_i(T_B\mathbf1)_i$. Nonnegativity yields
\[
 \sum_{B\in\Sigma^m}\psi(B)
 \le\sum_i\sum_{B\in\Sigma^m}(T_B\mathbf1)_i
 =\sum_i(T^m\mathbf1)_i\le N.
\]
If $w_L=B_1\cdots B_K$, submultiplicativity, $\|\alpha\|_1=1$,
and $\|\tau\|_\infty\le1$ imply
\[
 q(w_L)\le\prod_{r=1}^K\psi(B_r)
 \le\left(\frac NK\right)^K.
\]
The last inequality is arithmetic--geometric mean; it remains valid
when a factor vanishes. Thus
$-\log q(w_L)\ge L\log d-(L/m)\log N$.
Using $q(\varepsilon)\le1$ gives
\[
 D_{\rm KL}(p_L\|q)
 \ge-h_2(\delta_L)-\delta_L\log q(w_L),
\]
which proves the lower bound, including the case of infinite KL.

For the upper bound use the following two-state automaton:
\[
 \alpha=(1,0),\qquad
 \tau=\begin{pmatrix}1-M_0/L\\1/L\end{pmatrix},\qquad
 T_a=\begin{pmatrix}
 0&M_0/(Ld)\\0&(1-1/L)/d
 \end{pmatrix}\quad(a\in\Sigma).
\]
It is proper, has mean length $(M_0/L)L=M_0$, and satisfies
\[
 q_L(\varepsilon)=1-M_0/L,\qquad
 q_L(w)=\frac{M_0}{L^2}(1-1/L)^{L-1}d^{-L}
 \quad(|w|=L).
\]
Its divergence is exactly the upper expression in
\eqref{II:U:eq:kl-fixed-sandwich}. The two residual terms in that expression
vanish, as do $m^{-1}\log N$ and $h_2(\delta_L)$ for fixed $N$.

Each letter occurs $md^{m-1}=L/d$ times in $w_L$. Hence
Proposition~\ref{II:U:prop:iprojection} gives
$\mathcal K_1(p_L)=g(M_0)+M_0\log d-h_2(\delta_L)$.
Finally, a deterministic path of length $L$, with initial stop
probability $1-\delta_L$, initial continuation probability $\delta_L$,
and terminal stop probability one, realizes $p_L$ with $L+1$ states.
\end{proof}
